\documentclass[11pt]{article}
% =========================================================
%  學生講義 共用前言（以 XeLaTeX 編譯）
%  需要套件：xeCJK, tcolorbox（其餘多在 BasicTeX 內）
% =========================================================
\usepackage[a4paper,margin=2.3cm]{geometry}
\usepackage{amsmath,amssymb}
\usepackage{xcolor}
\usepackage{graphicx}

% ---- 中文字型（macOS 系統字型）----
\usepackage{xeCJK}
\setCJKmainfont{Songti TC}      % 內文：宋體
\setCJKsansfont{Heiti TC}       % 標題/強調：黑體
\xeCJKsetup{AutoFakeBold=2.2, AutoFakeSlant=0.2}

% ---- 版面 ----
\setlength{\parindent}{0pt}
\setlength{\parskip}{0.55em}
\linespread{1.15}
\renewcommand{\baselinestretch}{1.15}

% ---- 顏色 ----
\definecolor{cBlue}{HTML}{1f6feb}
\definecolor{cGray}{HTML}{6b7280}
\definecolor{cGrayBg}{HTML}{f3f4f6}
\definecolor{cPurp}{HTML}{7a3ff2}
\definecolor{cPurpBg}{HTML}{f5f1ff}

% ---- 標註框 ----
\usepackage[most]{tcolorbox}

% 就地補充新概念（灰底）：\begin{補充框}{標題}
\newtcolorbox{補充框}[1]{
  enhanced, breakable, arc=2pt, boxrule=0.6pt,
  colback=cGrayBg, colframe=cGray,
  left=9pt,right=9pt,top=7pt,bottom=7pt,
  before upper={\textbf{\sffamily #1}\par\vspace{3pt}},
}

% 延伸 / 開放問題（紫色左邊條）：\begin{延伸框}{標題}
\newtcolorbox{延伸框}[1]{
  enhanced, breakable, arc=2pt, boxrule=0pt,
  colback=cPurpBg, colframe=cPurpBg,
  borderline west={3pt}{0pt}{cPurp},
  left=10pt,right=9pt,top=7pt,bottom=7pt,
  before upper={\textbf{\sffamily 延伸閱讀｜#1}\par\vspace{3pt}},
}

% ---- 圖：置中、底下小字說明 ----
% \fig{檔名}{寬度}{說明}
\newcommand{\fig}[3]{%
  \begin{center}
  \includegraphics[width=#2\linewidth]{#1}\\[2pt]
  {\small\color{cGray} #3}
  \end{center}}

% ---- 章節標題用黑體 ----
\newcommand{\h}[1]{\sffamily #1}


\begin{document}

\begin{center}
{\sffamily\Huge\bfseries 必物-1 科學態度與方法}\\[3pt]
{\sffamily\large 普通高中 物理（必修）・學生講義}
\end{center}
\vspace{0.6em}

在開始計算任何物理之前，先弄清楚兩件事：\textbf{物理學是怎麼運作的}（如何從觀察走到理論、理論為什麼可以被推翻），以及\textbf{如何把世界量出來}（單位、因次、有效數字、數量級）。本章建立的是整年物理課所共用的工具與態度，後續每一章都會用到。

\medskip
本章主線：

\begin{center}
科學方法與可否證性 $\rightarrow$ 模型與近似 $\rightarrow$ 物理量與單位（SI）$\rightarrow$ 因次分析 $\rightarrow$ 測量、有效數字與不確定度 $\rightarrow$ 數量級與大自然的尺度
\end{center}

\section*{\sffamily 1-1\quad 科學的態度與方法}

\subsection*{\sffamily A.\ 物理學是「不斷被修正的最佳解釋」}

科學的核心不只是答案，而是\textbf{方法}。物理學處理的問題是：面對一個自然現象，如何提出解釋，又如何確認這個解釋值得相信。一個典型的流程是：

\begin{center}
觀察現象 $\rightarrow$ 提出假說（一個可檢驗的猜想）$\rightarrow$ 由假說推論出一個尚未觀察到的預測 $\rightarrow$ 設計實驗檢驗這個預測。
\end{center}

若預測符合，假說\textbf{暫時存活}（但不算「證明完畢」）；若預測不符，假說被推翻或須修正。其中有兩個關鍵：

\begin{enumerate}\setlength{\itemsep}{2pt}
\item \textbf{科學的進展主要靠「預測」，而不只是「解釋」}。任何結果發生後都能事後編出一套說法；較有力的是「理論預言接下來會發生 X」，而 X 確實發生。牛頓力學能算出哈雷彗星回歸的時間，廣義相對論預言星光經過太陽附近會被偏折——這種「先預測、後驗證」是科學的主要依據。
\item \textbf{科學結論是「暫時的」}。即使一個理論通過了大量檢驗，也只能說「目前尚未被推翻」，不能說「永遠正確」。牛頓力學沿用了兩百年，最後在高速與微觀尺度被相對論、量子力學修正。但牛頓力學並未「全錯」，它在其適用範圍內仍是極好的近似。
\end{enumerate}

\subsection*{\sffamily B.\ 可否證性：分辨科學與非科學的判準}

哲學家 Karl Popper 提出一個判準：一個主張要算科學，必須\textbf{原則上存在某種觀測結果，可以證明它是錯的}。這個性質稱為\textbf{可否證性（falsifiability）}。

\begin{itemize}\setlength{\itemsep}{1pt}
\item 「所有金屬受熱都會膨脹」：可否證——只要找到一塊受熱卻收縮的金屬，就能推翻它。這是科學陳述。
\item 「發生什麼都是注定的」：不可否證——無論發生什麼都能說「這就是安排」，沒有任何觀測能反駁它。這不屬於科學的範疇（不代表它沒有價值）。
\end{itemize}

一個理論若能做出大膽、具體、容易被否證的預測，它的科學含量越高。\textbf{能被推翻卻一直未被推翻}，正是信任一個理論的理由。

\begin{補充框}{什麼是可否證性？為什麼「不能被推翻」反而是缺點？}
可否證性的重點不在「對不對」，而在\textbf{有沒有可能錯給你看}。

\begin{itemize}\setlength{\itemsep}{1pt}
\item \textbf{科學陳述}「太陽明天從東邊升起」：只要哪天從西邊升起，它就被推翻。它把「怎樣算錯」講清楚了，正因為一直沒被推翻，才值得信任。
\item \textbf{不可否證陳述}「水星逆行使你諸事不順」：順利時可說「你命格強」，不順時可說「水星逆行」。無論結果如何都能自圓其說，於是沒有任何觀測能反駁它。
\end{itemize}

可否證性可類比為「敢不敢打賭，並先講好輸的條件」：能訂出「怎樣算我輸」的主張是可檢驗的；拒絕訂出輸贏條件的主張則無法檢驗。

算命、星座之所以常「感覺很準」，是因為用語\textbf{模糊到對誰都成立}（例如「最近會有貴人，也要小心小人」）。話講得越含糊，越多人覺得「在說我」，這稱為\textbf{巴納姆效應（Barnum effect）}。

\textbf{注意：}
\begin{itemize}\setlength{\itemsep}{1pt}
\item 可否證\textbf{不等於}已經被否證。它是「有可能被推翻」的資格，不是「已經被推翻」。
\item 不可否證\textbf{不等於}沒有價值。愛情、美學、道德、宗教信念都不可否證，只是不屬於科學的範疇。可否證性劃定的是「科學 vs.\ 非科學」，不是「有意義 vs.\ 沒意義」。
\item 科學無法證明一個理論「絕對為真」，只能不斷嘗試推翻它；通過越多嚴苛檢驗，越值得信任，但始終保留「將來可能被修正」的空間。
\end{itemize}
\end{補充框}

\subsection*{\sffamily C.\ 模型與近似：策略性的簡化}

物理學通常不直接處理「完整的真實世界」，而是處理\textbf{模型（model）}：抓住主要因素、暫時忽略次要因素的簡化版本。

\begin{itemize}\setlength{\itemsep}{1pt}
\item 計算棒球飛行：先把球當成\textbf{質點}（沒有大小、不旋轉），忽略空氣阻力；算得夠準後，再視需要把阻力、旋轉逐項加回。
\item 計算行星繞日：先把太陽與行星都當質點、軌道當正圓。
\end{itemize}

「忽略空氣阻力」不是偷懶，而是一種\textbf{策略性的近似}：先抓住最重要的物理，得到一個能計算、能理解的答案，再逐步逼近真實。判斷「什麼可以忽略」本身就是一種物理能力。

\textbf{注意：}
\begin{itemize}\setlength{\itemsep}{1pt}
\item 「模型就是不準」並不正確。模型是刻意取捨；好的模型在其適用範圍內既準確又好用。關鍵是\textbf{知道自己忽略了什麼、什麼時候不能再忽略}。
\item 「理論被修正等於之前的科學家錯了」並不必然。牛頓力學在日常尺度仍精準到足以將太空船送上月球；它是被「擴充適用範圍」，不是被「證明全錯」。
\end{itemize}

\subsection*{\sffamily D.\ 物理學研究的尺度範圍}

物理學試圖用\textbf{同一套基本定律}解釋從原子核（約 $10^{-15}$ m）到可觀測宇宙（約 $10^{26}$ m）的現象，橫跨超過 40 個數量級。這也是為什麼接下來要先學「如何有系統地描述如此懸殊的大小」——也就是 1-2 的單位與 1-4 的數量級。

\section*{\sffamily 1-2\quad 物理量與單位}

\subsection*{\sffamily A.\ 為什麼需要單位？基本量與導出量}

說「桌子長 3」沒有意義——3 個什麼？單位就是\textbf{約定好的比較基準}。報告一個物理量，必須同時給出「數值＋單位」。物理量分兩種：

\begin{itemize}\setlength{\itemsep}{1pt}
\item \textbf{基本量（base quantity）}：選定為「不靠其他量定義」的少數幾個量，其單位由標準直接規定。
\item \textbf{導出量（derived quantity）}：由基本量經乘除組合而來，例如速度＝長度÷時間、力＝質量×長度÷時間$^2$。
\end{itemize}

哪些量算「基本」其實是一種\textbf{人為的選擇}（也可以選另一套），但全世界統一採用一套，溝通才不會混亂。這套就是國際單位制。

\subsection*{\sffamily B.\ 國際單位制（SI）的七個基本量}

\textbf{國際單位制（SI，法文 Système International d'Unités）}選定七個基本量，各配一個基本單位。

\fig{m1-si-units}{0.70}{圖 1-1：SI 的七個基本量與基本單位。高中力學主要用到前三個——公尺、公斤、秒（合稱 MKS 制）。}

所有導出單位都由這七個基本單位組合而成，例如：
$$[\text{速度}]=\frac{\text{m}}{\text{s}},\qquad
[\text{力}]=\text{kg}\cdot\frac{\text{m}}{\text{s}^2}\equiv\text{N（牛頓）},\qquad
[\text{能量}]=\text{kg}\cdot\frac{\text{m}^2}{\text{s}^2}\equiv\text{J（焦耳）}.$$
牛頓（N）、焦耳（J）這些名稱是為了方便而取的；拆開來都是基本單位的組合。

\begin{補充框}{為什麼全世界要用同一套單位？「1 公斤」現在怎麼定義？}
\textbf{SI（國際單位制）}是全世界公定的測量基準，使任何人、任何國家量出的「1 公尺」「1 公斤」都指向\textbf{同一個量}，科學與貿易才能溝通。

\textbf{為什麼要統一？}\ 可類比為食譜：若一人說的「一杯」是 200 mL、另一人是 350 mL，照同一份食譜做菜結果就不同。科學更為嚴格：1999 年 NASA 的火星氣候探測者號，因一組工程師用英制磅、另一組用公制牛頓而未換算，探測器墜毀，損失上億美元。

\textbf{為什麼從「實物」改成「常數」？}\ 舊的「公斤」是 1889 年製作的一塊鉑銥合金圓柱，存放於巴黎。但它與各國複製品相比，一個多世紀來質量漂移了約 $50\ \mu\text{g}$；而且它只有一個，可能損壞或遺失。物理常數則處處相同、不隨時間改變：普朗克常數 $h$、光速 $c$ 在任何地點、任何時間都是同一個值。把單位「綁」在常數上，標準便不再漂移，且任何具備足夠儀器的實驗室都能自行重現，不必前往巴黎比對。

\textbf{現在怎麼定義（原理）}：把某個物理常數的數值「釘死」（固定為一個不再更動的數），反過來定義單位。以\textbf{公尺}為例：先把光速 $c$ 釘死為 $299\,792\,458\ \text{m/s}$，於是 1 公尺就被反過來定義成「光在 $1/299\,792\,458$ 秒走的距離」。常數一旦釘死，單位的大小也就跟著定下來，而且在任何地點、任何時間都能重現。

\fig{m1-redefine}{0.98}{圖 1-2：2019 年 SI 重新定義的核心機制。舊做法靠會漂移的實物標準（公斤原器），新做法把物理常數的數值釘死、反過來定出單位（以光速 $c$ 定出公尺為招牌例），標準便不再漂移、處處可重現。}

\textbf{2019 年起七個基本單位全部改用物理常數定義}，每個單位各對應一個被釘死的常數：
\begin{center}
\renewcommand{\arraystretch}{1.3}
\begin{tabular}{l l}
\hline
\textbf{基本單位} & \textbf{釘死的物理常數} \\
\hline
秒（s） & 銫-133 原子躍遷頻率 $\Delta\nu_{\text{Cs}}$ \\
公尺（m） & 光速 $c$ \\
公斤（kg） & 普朗克常數（Planck constant）$h$ \\
安培（A） & 基本電荷（elementary charge）$e$ \\
克耳文（K） & 波茲曼常數（Boltzmann constant）$k_{\text{B}}$ \\
莫耳（mol） & 亞佛加厥常數（Avogadro constant）$N_{\text{A}}$ \\
燭光（cd） & 發光效率（luminous efficacy）$K_{\text{cd}}$ \\
\hline
\end{tabular}
\end{center}

以\textbf{公斤}為例：把普朗克常數 $h$ 釘死為 $6.626\,070\,15\times10^{-34}\ \text{J·s}$；由於 $h$ 的單位含 kg，固定 $h$（再配合已定義的公尺、秒）即等於定義了公斤。重新定義時刻意讓新舊「1 公斤」在當時的最高精度下一致，因此日常數值完全無感，改變的是定義的根基，而非數值的大小。

\textbf{高中的重點}：記住招牌例子「\textbf{公斤}改用\textbf{普朗克常數} $h$ 定義」（以及「公尺對應光速 $c$」），其餘對照了解即可。
\end{補充框}

\subsection*{\sffamily C.\ 因次：只看「量的種類」}

把一個物理量「是長度、是時間、還是長度÷時間」這種\textbf{種類}抽出來，就是它的\textbf{因次（dimension）}。用方括號表示，力學三個基本因次記作：
$$[\text{長度}]=\mathsf{L},\qquad [\text{質量}]=\mathsf{M},\qquad [\text{時間}]=\mathsf{T}.$$
任何力學量的因次都能寫成 $\mathsf{M}^a\mathsf{L}^b\mathsf{T}^c$，例如：
$$[\text{速度}]=\mathsf{L}\,\mathsf{T}^{-1},\quad
[\text{加速度}]=\mathsf{L}\,\mathsf{T}^{-2},\quad
[\text{力}]=\mathsf{M}\,\mathsf{L}\,\mathsf{T}^{-2},\quad
[\text{能量}]=\mathsf{M}\,\mathsf{L}^2\,\mathsf{T}^{-2}.$$

\textbf{因次法則}：等號兩邊、相加減的每一項，因次必須完全相同。你可以把長度加長度，但不能把長度加時間。此外，三角函數、指數、對數的「輸入」必須是\textbf{無因次}（純數）的。

這條規則是\textbf{檢查公式的工具}：計算到一半若發現等號兩邊因次不符，公式必定抄錯或推錯。

\begin{補充框}{因次分析是什麼？為什麼用單位就能檢查公式？}
\textbf{因次分析（dimensional analysis）}的根本道理是一句話：\textbf{物理定律不會因為你選用什麼單位而改變}。太陽的質量不會因為用公斤量還是用磅量而不同，那只是同一件事的兩種記法。一條真正的物理等式，必須在換任何單位後依然成立，這個要求就逼出「等號兩邊因次必須相同」。

\textbf{為什麼兩邊因次要一樣？}\ 等號可類比為一座天平：兩邊各放「3 顆蘋果」可以比較數量；但「3 顆蘋果＝5 秒鐘」無法比較，因為種類不同，連「相不相等」都沒有意義。物理等式也一樣：左邊是「時間」，右邊就不能是「長度」。因此「兩邊因次相符」是公式正確的\textbf{必要條件}。

\textbf{因次分析能做兩件事：}
\begin{itemize}\setlength{\itemsep}{1pt}
\item \textbf{抓錯}：計算過程中檢查兩邊因次，不符就回頭找錯。
\item \textbf{定形式}：假設答案是相關量的乘冪組合，用因次定出各指數，可推出公式的「骨架」。
\end{itemize}

\textbf{注意：}因次分析\textbf{定不出無因次的純數}（如 $2\pi$、$\tfrac12$）。例如 $T=2\pi\sqrt{L/g}$ 和 $T=99\sqrt{L/g}$ 因次完全相同。所以因次相符只是必要條件，不是充分條件：因次相符\textbf{不保證}公式正確。此外，因次（種類，如長度）與單位（該種類的某個刻度，如公尺、公分、英尺）是兩回事，不可混為一談。
\end{補充框}

\section*{\sffamily 1-3\quad 測量、有效數字與不確定度}

\subsection*{\sffamily A.\ 任何測量都有誤差}

以尺量桌子，最終只能讀到「介於某兩刻度之間」，最後一位是估出來的。換一把更精密的尺，估的位數往後挪，但「最後一位仍是估的」這件事不會消失。測量的本質是\textbf{逼近}，而非「捕捉到真值」，所以每個測量值都帶有一段\textbf{不確定度（uncertainty）}。「測量必有誤差」不只是「儀器不夠好」的技術問題，背後還牽涉更根本的物理，詳見 1-3 末的延伸框。

\subsection*{\sffamily B.\ 準確度與精密度}

準確度與精密度是兩件不同的事，常被混用。以射靶來理解最清楚。

\fig{m1-accuracy}{0.95}{圖 1-3：以彈著點分布說明準確度與精密度。由左至右為：低準確度低精密度、高準確度低精密度、低準確度高精密度、高準確度高精密度。靶心代表真值。}

\begin{itemize}\setlength{\itemsep}{1pt}
\item \textbf{準確度（accuracy）}：測量值\textbf{離真值多近}（彈著點離靶心多近）。準確度差代表有\textbf{系統誤差}（如尺零點未對齊、儀器未歸零），整批一起偏移。
\item \textbf{精密度（precision）}：重複測量\textbf{彼此多集中}（彈著點之間多密）。精密度差代表\textbf{隨機誤差}大，每次都有抖動。
\end{itemize}

兩者互相獨立：可以「很集中但整批偏掉」（精密但不準），也可以「鬆散卻平均落在靶心」（準但不精密）。

\begin{center}
\renewcommand{\arraystretch}{1.3}
\begin{tabular}{p{0.20\linewidth} p{0.33\linewidth} p{0.33\linewidth}}
\hline
 & \textbf{系統誤差} & \textbf{隨機誤差} \\
\hline
表現 & 整批往同方向偏 & 每次忽高忽低 \\
影響 & 準確度 & 精密度 \\
如何減少 & 校正儀器、改良方法 & 多測幾次取平均 \\
\hline
\end{tabular}
\end{center}

\textbf{注意：}「多測幾次取平均」只能壓低\textbf{隨機}誤差，對\textbf{系統}誤差無效。一把短了 1 mm 的尺，量一百次取平均仍然短 1 mm。

\subsection*{\sffamily C.\ 有效數字}

一個測量值中「真正攜帶資訊」的數字，就是\textbf{有效數字（significant figures）}。多寫的位數等於假裝量得比實際更準，並不誠實。

\textbf{如何判斷哪些位數算「有效」：}
\begin{enumerate}\setlength{\itemsep}{1pt}
\item 非零數字一律有效：$3.14$ 有 3 位。
\item 數字之間的 0 有效：$2.005$ 有 4 位。
\item 小數點後、末尾的 0 有效（表示量到了那一位）：$1.20$ 有 3 位，比 $1.2$（2 位）更精確。
\item 純粹定位的前導 0 不算：$0.0042$ 只有 2 位。
\item 像 $1500$ 這種結尾 0 有歧義；改用科學記號即無爭議：$1.5\times10^3$（2 位）、$1.50\times10^3$（3 位）。
\end{enumerate}

\textbf{運算規則：}
\begin{itemize}\setlength{\itemsep}{1pt}
\item \textbf{乘除}：結果的有效位數取\textbf{最少}的那個輸入。例：$2.3\times1.58=3.634\to3.6$（2 位，因 $2.3$ 只有 2 位）。
\item \textbf{加減}：結果對齊到「小數位最少」的那一項（看小數點後位數，不是看有效位數）。例：$12.11+1.1=13.21\to13.2$（對齊到小數第一位）。
\item \textbf{先算完、最後一步才修約}；中間過程多留 1\,$\sim$\,2 位，以免誤差累積。
\end{itemize}

\textbf{注意：}
\begin{itemize}\setlength{\itemsep}{1pt}
\item 「計算機顯示越多位就越精確」並不正確。輸入只有 2、3 位，輸出就不應超過。把 $\pi r^2$ 算出的一長串數字全部抄上去，是假精確。
\item 加減看「小數位」、乘除看「有效位數」，兩條規則不同，不可混用。
\item $0.0500$ 有 3 位有效（前導 0 不算，末尾兩個 0 算）。
\end{itemize}

\subsection*{\sffamily D.\ 科學記號與數量級估計}

懸殊的大小用\textbf{科學記號} $a\times10^{n}$（$1\le |a|<10$）表示最清楚，也使有效位數一目了然。只看 10 的次方（不在意前面的係數），就是一個量的\textbf{數量級（order of magnitude）}。

物理上常用的習慣是：\textbf{先估數量級，再決定要不要細算}。這種「用粗略假設拼出合理估值」的方法稱為\textbf{費米估計（Fermi estimation）}，得名自能徒手估出原子彈當量的 Enrico Fermi。它的精神不是算得精確，而是先判斷答案大約落在 10、100、1000 還是更高的量級，藉此立刻判斷一個數字是否合理。例如有人聲稱「一秒鐘心跳一萬下」：正常心跳約每秒 1\,$\sim$\,2 下（數量級 $10^0$），一萬（$10^4$）比它大了四個數量級，物理上不可能。學會估數量級，等於有了一個隨身的合理性檢查工具。

\begin{延伸框}{有沒有「絕對精確」的測量？}
課本把不確定度當成「儀器不夠好」的技術問題，言下之意是「只要儀器無限精密，就能量到絕對精確」。但這只停在實務技術層次，沒有回答更根本的問題：誤差只是技術限制，還是自然本身就設了下限？不確定度其實來自三個層層更深的來源。

\textbf{第一層：讀數與儀器（技術性，可改善）}。任何刻度尺，最後一位都是估出來的；換更細的尺只是把估的位數往後挪，估的動作本身消不掉。儀器也有解析度極限。這層可靠更好的工具持續改善，但永遠到不了零。

\textbf{第二層：被測對象與環境的擾動（統計性，可壓低）}。要量的東西本身常在波動：空氣分子不停碰撞、溫度微幅起伏、桌子有振動。多測取平均能把這種隨機抖動壓很低（測 $N$ 次，隨機誤差約縮為 $1/\sqrt{N}$），但壓不到零。

\textbf{第三層：量子力學的下限（原理性，無法消除）}。\textbf{海森堡測不準原理（uncertainty principle）}指出：某些成對的物理量無法同時被無限精確地確定，這不是儀器問題，而是自然的規定。最著名的是位置 $x$ 與動量 $p$：
$$\Delta x\,\Delta p \gtrsim \frac{\hbar}{2}.$$
把位置量得越準（$\Delta x$ 越小），動量就必然越模糊（$\Delta p$ 越大），反之亦然，乘積有一個不可逾越的下限。可類比為：要「看見」一顆電子的位置，得用光去照它，但光一打上去就把電子撞動了。不過真正的測不準\textbf{比這更根本}——它不是「測量手段笨拙」，而是電子本身就沒有同時確定的位置與動量。（這條線到 \textbf{必物-6 量子現象} 會正式展開。）

因此：\textbf{不存在「絕對精確」的測量}。前兩層使你「實務上」永遠差一點，第三層使你「原理上」連想都不能想。

\textbf{哪些是定論、哪些仍開放：}「測量必有下限、不存在零不確定度的測量」在標準量子力學框架內是\textbf{確定的結論}，已被大量實驗證實。但測不準原理「為什麼」成立——也就是量子力學對「測量」的根本詮釋——\textbf{目前尚無定論}。微觀粒子在被測量之前到底有沒有確定的位置，不同的量子詮釋給的答案不同，這牽涉「量子測量問題」，是物理學至今尚無共識的開放領域。

\textbf{注意：}不確定度不是「出錯」或「粗心」，而是對「我們知道得多準」的誠實標示，是科學嚴謹的表現。一個沒有附上不確定度的測量值，反而不夠專業。
\end{延伸框}

\section*{\sffamily 1-4\quad 大自然的尺度}

把本章學到的單位與數量級用來描述整個宇宙，可作為本章的收尾。下圖把長度從原子核排到可觀測宇宙，用的是\textbf{對數軸}——每往右一格，大小就乘 10。

\fig{m1-scales}{0.95}{圖 1-4：大自然的長度尺度（對數軸，每格代表 10 倍）。橫跨從原子核（約 $10^{-15}$ m）到可觀測宇宙（約 $10^{26}$ m）。}

\textbf{幾個務必記住的長度數量級：}
\begin{center}
\renewcommand{\arraystretch}{1.2}
\begin{tabular}{l l @{\hskip 2.5em} l l}
\hline
對象 & 數量級 (m) & 對象 & 數量級 (m) \\
\hline
原子核 & $10^{-15}$ & 地球直徑 & $10^{7}$ \\
原子 & $10^{-10}$ & 地球–太陽距離 & $10^{11}$ \\
病毒 & $10^{-7}$ & 銀河系直徑 & $10^{21}$ \\
人 & $10^{0}$ & 可觀測宇宙 & $10^{26}$ \\
\hline
\end{tabular}
\end{center}

兩個最該記住的物理尺度是：\textbf{原子約 $10^{-10}$ m、原子核約 $10^{-15}$ m}。兩者相差 5 個數量級（十萬倍）。若把原子放大到一座體育場，原子核大約只有場中央的一顆綠豆——原子內部幾乎是空的。這個事實到\textbf{必物-3 物質的組成與交互作用}會再用到。

時間與質量也各有自己的尺度跨度（參考）：
\begin{itemize}\setlength{\itemsep}{1pt}
\item \textbf{時間}：光穿過原子核約 $10^{-23}$ s；人心跳約 $10^{0}$ s；人壽命約 $10^{9}$ s；宇宙年齡約 $10^{17}$ s（約 138 億年）。
\item \textbf{質量}：電子約 $10^{-30}$ kg；一滴水約 $10^{-4}$ kg；人約 $10^{2}$ kg；地球約 $10^{24}$ kg；太陽約 $10^{30}$ kg。
\end{itemize}

\textbf{注意：}對數軸上「等距離＝等倍率」，不是「等差距」。從 $10^{0}$ 到 $10^{5}$ 與從 $10^{20}$ 到 $10^{25}$ 在圖上一樣長，因為兩者都是「乘了 $10^5$ 倍」。第一次看對數軸時容易誤以為間距代表加減。

物理學的目標，是用同一套定律描述從原子核到可觀測宇宙、橫跨約 $10^{41}$ 倍的世界（$10^{26}/10^{-15}=10^{41}$）。這與 1-1 所說「用同一套基本定律解釋一切」前後呼應。

\section*{\sffamily 1-5\quad 本章重點整理}

\begin{補充框}{一頁複習（關鍵結論）}
\begin{itemize}\setlength{\itemsep}{2pt}
\item \textbf{科學方法}：觀察 $\rightarrow$ 假說 $\rightarrow$ 預測 $\rightarrow$ 實驗檢驗 $\rightarrow$ 修正；結論永遠「暫時為真」。
\item \textbf{可否證性}：能被某種觀測推翻才算科學；「怎樣都對」反而不科學。可否證是「有可能被推翻」的資格，不等於已被推翻。
\item \textbf{模型與近似}：刻意簡化、抓住主要因素（質點、忽略阻力）；關鍵是知道忽略了什麼、何時不能再忽略。
\item \textbf{SI 七個基本量}：長度 m、質量 kg、時間 s、電流 A、溫度 K、物量 mol、發光強度 cd；力學常用 MKS。導出量＝基本量的乘除組合（$\text{N}=\text{kg}\cdot\text{m/s}^2$）。2019 年起七個基本單位全改用物理常數定義。
\item \textbf{因次}：相加減與等號兩邊因次須一致；$[v]=\mathsf{L}\mathsf{T}^{-1}$、$[a]=\mathsf{L}\mathsf{T}^{-2}$、$[F]=\mathsf{M}\mathsf{L}\mathsf{T}^{-2}$。可抓錯、定形式，但定不出無因次純數；因次相符是必要條件，非充分條件。
\item \textbf{準確度與精密度}：準＝離真值近（壓系統誤差靠校正）；密＝彼此集中（壓隨機誤差靠多測取平均）。多測取平均對系統誤差無效。
\item \textbf{有效數字}：乘除取最少有效位數、加減對齊最少小數位；最後一步才修約。
\item \textbf{數量級與費米估計}：先估 10 的次方判斷合理性；科學記號 $a\times10^n$。
\item \textbf{尺度}：原子 $10^{-10}$ m、原子核 $10^{-15}$ m、可觀測宇宙 $10^{26}$ m；全宇宙橫跨約 $10^{41}$ 倍。對數軸上等距離代表等倍率。
\item \textbf{不確定度}：任何測量都有不確定度，「絕對精確」的測量不存在；最底層的量子下限（測不準原理）無法消除。
\end{itemize}
\end{補充框}

\end{document}
